%% defense_mechanism.tex

\section{The \reauthgate{} Design Pattern}
\label{sec:defense}

We propose a single OS-principled design pattern---the \reauthgate{}---that ports the classical OS exception-handler privilege invariant to the LLM agent recovery path. The pattern is implemented in the \system{} codebase and validated in Section~\ref{sec:evaluation}. It is framework-agnostic in principle: it depends only on the existence of a recovery orchestrator that mediates side-effecting recovery actions and a permission checker that gates the forward execution path, not on any specific agent architecture.

\subsection{Motivation: Centralized vs.\ Decentralized Validation}
\label{sec:centralize-motivation}

The root cause of CWE-862 at the recovery sink (\S\ref{sec:threat-model}) is the \emph{decentralized} validation of recovery actions. In the baseline design, each recovery strategy is responsible for its own permission checks: \code{TopologyMutationStrategy} is expected to call \code{PermissionChecker} before replacing a role, \code{ModelFallbackStrategy} before switching models, and so on. In practice this leads to two failure modes that are well-known in the cross-cutting concern literature~\cite{aspectj}:

\begin{enumerate}[leftmargin=*]
    \item \textbf{Gaps.} A strategy author who forgets or mis-scopes a check (as in the baseline \code{parseAndValidate(validate=false)} call) leaves a silent hole, because no other component re-checks the action. The baseline vulnerability in \system{} is exactly such a gap: the \code{validate=false} flag was a deliberate experimental switch, but in a production codebase, an equivalent gap could arise from a forgotten check in a newly added strategy.
    \item \textbf{Inconsistency.} Different strategies implement slightly different policies, so the effective privilege boundary depends on which strategy the orchestrator happens to dispatch. A role replacement via \code{TopologyMutationStrategy} might be checked, while a role replacement via a future \code{RolePromotionStrategy} might not be.
\end{enumerate}

This is a classic instance of a \emph{cross-cutting concern} in the software-engineering sense: permission enforcement is an aspect that \emph{every} side-effecting recovery action needs, but it is orthogonal to the strategy's core logic (failure diagnosis and recovery planning). Scattering the concern across strategies couples diagnosis logic to security logic, duplicates enforcement code, and---most importantly---makes the security property \emph{non-local}: there is no single place where one can verify that ``no recovery action bypasses permission checks.''

The principled fix is to centralize the concern at a single choke point. We do this by (i)~extending the recovery result type with a \code{requiresReauthorization} flag, (ii)~deferring the side effect as a callback stored in the recovery context rather than executing it inline, and (iii)~intercepting every flagged result at the \code{RecoveryOrchestrator}---the single component through which every recovery action must flow. The strategies declare \emph{that} an action needs reauthorization; the orchestrator enforces \emph{how}. This separation localizes the security property: the invariant ``no side effect executes without reauthorization'' is now verifiable by inspecting one method rather than fourteen strategies.

\begin{figure}[htbp]
\centering
\begin{tikzpicture}[
    node distance=0.7cm and 0.8cm,
    box/.style={draw, rounded corners=2pt, minimum height=0.7cm, minimum width=3.0cm, align=center, font=\footnotesize},
    strat/.style={box, fill=yellow!20, draw=orange!70},
    result/.style={box, fill=orange!20, draw=orange!70, font=\footnotesize\ttfamily},
    gate/.style={box, fill=yellow!30, draw=orange!80, minimum width=3.5cm, font=\footnotesize\bfseries},
    check/.style={box, fill=orange!25, draw=orange!70, minimum width=3.0cm},
    pass/.style={box, fill=green!15, draw=green!60!black},
    block/.style={box, fill=red!15, draw=red!70!black},
    note/.style={box, fill=gray!10, draw=gray!50, font=\tiny\itshape, minimum width=4.5cm},
    arr/.style={-Stealth, thick},
    parr/.style={-Stealth, thick, green!60!black},
    barr/.style={-Stealth, thick, red!70!black},
    lbl/.style={font=\tiny\itshape, gray!60!black, align=center}
]
% Top flow
\node[strat] (s) {Recovery Strategy\\\code{TopologyMutation}};
\node[result, below=of s] (r) {RecoveryResult\\reqReauth=true};
\node[gate, below=of r] (g) {\reauthgate{}\\\code{RecoveryOrchestrator}};
\node[check, below=of g] (c) {PermissionChecker\\$P_{new} \leq P_{old}$?};

% Branches
\node[pass, below left=0.8cm and 0.3cm of c] (exec) {Execute Side Effect\\resumeNode(newRole)};
\node[block, below right=0.8cm and 0.3cm of c] (deny) {Block\\RecoveryResult.failed};

\node[pass, below=0.6cm of exec] (ok) {Recovery OK\\(privilege preserved)};
\node[block, below=0.6cm of deny] (no) {Attack Stopped};

\draw[arr] (s) -- (r) node[lbl, midway, right] {declare};
\draw[arr] (r) -- (g) node[lbl, midway, right] {intercept};
\draw[arr] (g) -- (c) node[lbl, midway, right] {check};
\draw[parr] (c) -- (exec) node[lbl, midway, above] {pass};
\draw[barr] (c) -- (deny) node[lbl, midway, above] {fail};
\draw[parr] (exec) -- (ok);
\draw[barr] (deny) -- (no);

% Invariant note
\node[note, left=1.2cm of c] (inv) {\privinv{}:\\$P(\text{action}) \leq P(\text{failed exec})$\\(analog: OS exception handler\\never escalates privileges)};
\draw[dashed, gray!50] (inv) -- (c);
\end{tikzpicture}
\caption{The \reauthgate{} architecture. Validation is decoupled from strategy implementation: strategies declare the need for reauthorization via a flag, and a single orchestrator-level gate enforces the \privinv{} before any side effect executes. The invariant (left) constrains the checker.}
\label{fig:reauthgate-arch}
\end{figure}

\subsection{Principle: The \privinv{}}
\label{sec:invariant}

The \reauthgate{} enforces the \privinv{}: \emph{no recovery action may exceed the privilege of the failed execution it remedies}. Formally, let $P(a)$ denote the privilege level of recovery action $a$ and $P_{\text{fail}}$ the privilege of the failed execution. Then
\[
P(a) \leq P_{\text{fail}}.
\]
This is the agent-systems analog of the OS exception-handler privilege rule: an exception handler executes with the privileges of the faulting context, never higher~\cite{levine2003system,bovet2005understanding}. The same principle underlies the Confused Deputy fix~\cite{hardy1988confused}: the privileged deputy may not exercise authority not granted to the caller, even when the caller's request looks superficially legitimate.

For role replacement, the invariant specializes to: the new role's privileges must not exceed the original role's privileges. The \code{PermissionChecker} implements this via \code{checkRoleMutation(originalRole, newRole)}, which verifies that the new role is in the pre-approved downgrade set for the original role (e.g., \code{Python\_Coder} $\rightarrow$ \code{Code\_Reviewer} is allowed; \code{Code\_Reviewer} $\rightarrow$ \code{System\_Admin} is blocked). Combined with a role-existence whitelist---the role must be registered in the blueprint registry, not an arbitrary LLM-suggested string---the gate closes both the CWE-862 sink (the missing authorization check) and the Confused Deputy vector (the orchestrator no longer exercises role-replacement authority on the basis of attacker-influenced diagnostic content).

\subsection{Design: Declare at the Strategy, Enforce at the Orchestrator}
\label{sec:design-decomposition}

The pattern decomposes into three concrete code-level mechanisms, each grounded in established software-engineering practice for managing cross-cutting concerns~\cite{aspectj}.

\subsubsection{Mechanism 1: Result-Type Extension}
The \code{RecoveryResult} class is extended with a boolean \code{requiresReauthorization} flag (default: \code{false}). Strategies that produce side effects---role replacement, tool execution, model switching---set this flag to \code{true} via a factory method \code{okRequiringReauthorization(message, payload)}. Purely diagnostic strategies (e.g., reflection injection that only appends to the prompt context) leave the flag at its default. The flag is a \emph{declaration}, not a check: it tells the orchestrator ``this result needs reauthorization before its side effect executes,'' without specifying what the authorization check should be.

\subsubsection{Mechanism 2: Deferred Side Effects}
A side-effecting strategy does \emph{not} execute its side effect inline. Instead, it stores the side effect as a \emph{deferred callback} in the recovery context's metadata, keyed by \code{META\_PENDING\_SIDE\_EFFECT}. The callback is a \code{Function<RecoveryContext, RecoveryResult>}---a closure over the recovery context that, when invoked, performs the actual side effect (e.g., calling \code{resumeNode}). This deferral is critical: it ensures that the side effect is \emph{not} executed at strategy-return time, but is held in suspension until the orchestrator has had a chance to intercept and reauthorize.

\begin{lstlisting}[caption={Deferring the side effect in \code{TopologyMutationStrategy}; the strategy declares reauthorization need but does not execute the side effect.},label={lst:defer-side-effect}]
// In TopologyMutationStrategy.apply()
context.withMetadata(RecoveryOrchestrator.META_PENDING_SIDE_EFFECT,
    (RecoveryContext ctx) -> {
        boolean mutated = WorkflowEngine.instance()
            .resumeNode(nodeId, report, workflowId);
        return mutated
            ? RecoveryResult.ok("Topology mutation applied")
            : RecoveryResult.failed("resumeNode returned false");
    });
return RecoveryResult.okRequiringReauthorization(
    "Topology mutation pending reauth: proposed " + suggestedRole, null);
\end{lstlisting}

\subsubsection{Mechanism 3: Orchestrator-Level Interception}
The \code{RecoveryOrchestrator}'s main recovery loop intercepts every result with \code{requiresReauthorization=true} \emph{before} executing the deferred callback. The gate calls the existing \code{PermissionChecker} (the same one used on the forward path) to verify the \privinv{}. For role replacement, the gate consults \code{RoleReplacementValidator}, which performs two checks: (i)~the proposed role exists in the blueprint whitelist, and (ii)~the proposed role is not a privilege escalation from the failing agent's current role. If either check fails, the side effect is blocked and the orchestrator returns a failed recovery result; if both pass, the deferred callback executes (Listing~\ref{lst:reauthgate}).

\begin{lstlisting}[caption={\reauthgate{} in \code{RecoveryOrchestrator}: side effect deferred until reauth passes.},label={lst:reauthgate}]
// In RecoveryOrchestrator.dispatchRecovery()
if (result.requiresReauthorization()) {
    ReauthResult reauth = RecoveryReauthorizationGate.check(
        result, context, resolvePermissionChecker());
    if (!reauth.allowed()) {
        log.warn("[ReauthGate] blocked recovery action: {}", reauth.reason());
        return RecoveryResult.failed("denied by reauth gate: " + reauth.reason());
    }
    // reauth passed -> execute deferred side effect
    Function<RecoveryContext, RecoveryResult> sideEffect =
        context.metadata().get(META_PENDING_SIDE_EFFECT);
    result = sideEffect.apply(context);
}
\end{lstlisting}

This is a \emph{prevent} control, not a \emph{detect-and-alert} control: a rejected reauthorization blocks the side effect entirely, rather than merely logging it after the fact. The distinction matters because a privilege escalation that has already happened cannot be undone by an after-the-fact alert---the agent has already been replaced with a higher-privilege role.

\subsection{Why This Beats Per-Strategy Checks}
\label{sec:why-centralize-pattern}

Centralizing reauthorization at the orchestrator yields four concrete advantages over the decentralized baseline, each grounded in established software-engineering practice for managing cross-cutting concerns:

\begin{enumerate}[leftmargin=*]
    \item \textbf{Unified policy.} All security-critical recovery actions traverse the same gate, so the effective policy is identical regardless of which strategy produced the action. This eliminates the inconsistency and gap failure modes of decentralized validation (\S\ref{sec:centralize-motivation}).
    \item \textbf{Prevent, not detect.} Because side effects are deferred until reauthorization passes, an unauthorized action is blocked \emph{before} any state change occurs. A detect-and-alert design (e.g., auditing \code{resumeNode} calls after execution) cannot undo a privilege escalation that has already happened.
    \item \textbf{Separation of concerns.} Recovery strategies focus on failure diagnosis and recovery planning; the security aspect is factored out into a dedicated component. Strategy authors no longer need to be security experts, and security authors no longer need to understand every strategy's internals. This is the same separation that motivates aspect-oriented programming~\cite{aspectj} and middleware-based enforcement in classical distributed systems.
    \item \textbf{Open/closed for extension.} A newly added recovery strategy inherits protection automatically by setting \code{requiresReauthorization=true}; it does \emph{not} need to re-implement permission checks. The security property is thus \emph{open for extension} (new strategies are easy to add) but \emph{closed for modification} (the gate code does not change). This scales to the 14 strategies in \system{}'s orchestrator without per-strategy patches, and means that the security invariant cannot regress when a new strategy is added.
\end{enumerate}

The net effect is that the security property ``no recovery action exceeds the privilege of the failed execution'' becomes a \emph{local, inspectable} property of one orchestrator method, rather than an \emph{emergent} property of fourteen independently-authored strategies. This locality is itself an OS principle: classical kernels enforce privilege invariants at well-defined boundaries (syscall entry, exception dispatch), not inside each driver or subsystem.

\subsection{Layered Defense: Whitelist at the Strategy, Gate at the Orchestrator}
\label{sec:layered-defense}

The \system{} implementation actually employs the \reauthgate{} pattern in two layers, both of which enforce the same \privinv{} but at different points in the recovery pipeline:

\begin{enumerate}[leftmargin=*]
    \item \textbf{Layer 1 (strategy-level): \code{RoleReplacementValidator}.} The \code{TopologyMutationStrategy.parseAndValidate} method calls \code{RoleReplacementValidator.validate(currentRole, suggestedRole)} \emph{before} declaring the reauthorization need. The validator performs: (a)~existence whitelist (the suggested role must be in the blueprint registry), and (b)~non-escalation check (the suggested role's privileges must not exceed the current role's). A failure at Layer 1 short-circuits the strategy and returns a failed result without setting \code{requiresReauthorization}; no side effect is even deferred.
    \item \textbf{Layer 2 (orchestrator-level): \code{RecoveryReauthorizationGate}.} Even if Layer 1 passes, the orchestrator-level gate re-checks the invariant before executing the deferred side effect. This is defense-in-depth: a bug in Layer 1 (e.g., a future strategy that forgets to call the validator) is caught by Layer 2.
\end{enumerate}

The two layers are not redundant: Layer 1 is a fast-fail optimization that avoids the overhead of deferring a callback that will be rejected, while Layer 2 is the structural guarantee that closes the gap if a strategy bypasses or mis-implements Layer 1. The empirical validation (\S\ref{sec:evaluation}) reports results with both layers enabled; toggling Layer 1 off (the \code{validate=false} baseline) re-opens the vulnerability, confirming that Layer 2 alone suffices to block the attack but Layer 1 reduces wasted work.

\subsection{Framework-Agnostic Pattern Specification}
\label{sec:pattern-spec}

The \reauthgate{} pattern is specified independent of \system{}'s implementation. A framework porting the pattern needs:

\begin{enumerate}[leftmargin=*]
    \item A \emph{recovery orchestrator} that mediates every side-effecting recovery action (the single choke point). Most self-healing frameworks already have such a component~\cite{autogen,langgraph,reflexion}.
    \item A \emph{permission checker} that gates the forward execution path. The pattern reuses this checker, not a new one.
    \item A \emph{recovery result type} that strategies return, extensible with a boolean flag.
    \item A \emph{side-effect deferral mechanism}---a way to package a side effect as a callback that the orchestrator can invoke later. In Java this is a \code{Function<Context, Result>}; in Python it would be a closure; in TypeScript a \code{() => Result} function.
\end{enumerate}

The pattern does \emph{not} require: a specific agent architecture (graph-based or otherwise), a specific LLM, a specific role/permission model, or any modification to the forward execution path. The migration cost is concentrated in three sites: the result type (add a flag), the orchestrator loop (add an interception branch), and each side-effecting strategy (set the flag and defer the side effect). In \system{}, the full migration touched 14 strategy classes and one orchestrator method, with sub-millisecond runtime overhead and zero false alarms on benign recovery (\S\ref{sec:evaluation}).
